\section{Servidor de base de datos (169.58.62.99)}

\subsection{Archivos necesarios}
\pendiente{\texttt{podman-compose-bd.sh}, \texttt{database-scripts/init/00-init.sql} y \texttt{database-scripts/seeds/}}

\subsection{Descripción del script podman-compose-bd.sh}
\pendiente{comandos disponibles: up, reset, seed, firewall, firewall-boot, backup, restore, psql-*}

\subsection{Volúmenes y persistencia de datos}
\pendiente{dónde viven los datos de PostgreSQL y qué pasa al recrear el contenedor}

\subsection{Creación de defensoria-db (puerto 5432)}
\pendiente{comando, variables POSTGRES\_*, base \texttt{defensoria\_db}}

\subsection{Creación de historico-db (puerto 5433)}
\pendiente{base \texttt{defensoria\_historico\_db}; la carga de datos históricos queda a cargo de otra persona}

\subsection{Script de inicialización y esquema}
\pendiente{qué crea 00-init.sql y qué tablas genera Hibernate (\texttt{ddl-auto: update}) al arrancar los servicios}

\subsection{Carga de catálogos semilla}
\pendiente{01-dependencias, 02-chatbot, 03-personal\_test (este último solo para pruebas)}

\subsection{Firewall: acceso solo desde el Backend}
\pendiente{\texttt{./podman-compose-bd.sh firewall} y \texttt{firewall-boot}; los puertos 5432/5433 solo aceptan la IP 156.67.26.73}

\subsection{Verificación}
\begin{lstlisting}
# podman ps
# ./podman-compose-bd.sh psql-defensoria   # \dt
\end{lstlisting}
